% ─────────────────────────────────────────────────────────────────────────────
% Ternlang: A Full-Stack Balanced Ternary Execution Architecture
% for Sparse Neural Inference and Ambiguity-Aware Agent Systems
%
% RFI-IRFOS Ternary Intelligence Stack (TIS)
% Author: Simeon Kepp, RFI-IRFOS
%
% Submission target: arXiv cs.PL / cs.AR / cs.NE
% Format: IEEE two-column (IEEEtran) or ACM (acmart)
% ─────────────────────────────────────────────────────────────────────────────

\documentclass[journal,10pt,twocolumn]{IEEEtran}

\usepackage{amsmath,amssymb,amsthm}
\usepackage{listings}
\usepackage{xcolor}
\usepackage{booktabs}
\usepackage{graphicx}
\usepackage{hyperref}
\usepackage{microtype}
\usepackage{multirow}
\usepackage{array}

% ── Ternlang code style ───────────────────────────────────────────────────────
\definecolor{ternkw}{RGB}{0,164,140}
\definecolor{terncmt}{RGB}{120,120,150}
\definecolor{ternlit}{RGB}{220,80,80}
\definecolor{ternbg}{RGB}{18,18,30}
\definecolor{ternfg}{RGB}{220,220,230}

\lstdefinelanguage{ternlang}{
  morekeywords={fn,let,match,if,for,in,loop,break,return,
                agent,spawn,send,await,struct,trit,trittensor,agentref,
                truth,hold,conflict,consensus,invert,matmul,sparsity,cast},
  morecomment=[l]{//},
  morestring=[b]",
  sensitive=true
}

\lstset{
  language=ternlang,
  basicstyle=\ttfamily\footnotesize,
  keywordstyle=\color{ternkw}\bfseries,
  commentstyle=\color{terncmt}\itshape,
  stringstyle=\color{ternlit},
  showstringspaces=false,
  breaklines=true,
  frame=single,
  numbers=left,
  numberstyle=\tiny\color{terncmt},
  numbersep=4pt,
  xleftmargin=10pt
}

\newtheorem{definition}{Definition}
\newtheorem{theorem}{Theorem}
\newtheorem{proposition}{Proposition}

% ─────────────────────────────────────────────────────────────────────────────

\title{Ternlang: A Full-Stack Balanced Ternary Execution Architecture
       for Sparse Neural Inference and Ambiguity-Aware Agent Systems}

\author{Simeon~Kepp,~\IEEEmembership{RFI-IRFOS}%
\thanks{Correspondence: RFI-IRFOS Ternary Intelligence Stack.
Repository: \url{https://github.com/eriirfos-eng/ternary-intelligence-stack--tis-}}%
}

\begin{document}

\maketitle

% ─────────────────────────────────────────────────────────────────────────────
\begin{abstract}
We present \textbf{Ternlang}, the first complete software stack for
balanced ternary computing: a domain-specific language, bytecode compiler,
stack-based virtual machine (BET VM), hardware description language backend,
distributed actor runtime, and machine learning inference kernels—all unified
under a single coherent architecture. The foundational primitive is the
\emph{trit} $t \in \{-1, 0, +1\}$, where the value $0$ represents an
\emph{active neutral state} rather than absence, enabling three-valued logic
that is structurally superior to binary for ambiguity-aware reasoning.

The principal contribution is \textbf{TSPARSE\_MATMUL}: a first-class VM
opcode that elides multiply operations against zero-weighted (``hold'') trit
elements, surfacing at the instruction-set level a property that
BitNet-style~\cite{ma2024era} ternary quantization reveals in neural weights.
Empirical evaluation on quantized weight matrices of $512 \times 512$
elements yields $56.2\%$ sparsity and a $2.27\times$ reduction in multiply
operations versus dense execution; a Compressed Sparse Column kernel with
Rayon parallelism achieves up to $122\times$ speedup at $99\%$ sparsity—
without approximate arithmetic or hardware reconfiguration.

We further define the \textbf{Balanced Ternary Execution (BET) ISA}, a
formal 2-bit-encoded instruction set with 51 opcodes spanning arithmetic,
tensor operations, actor messaging, and control flow; synthesise it to
Verilog-2001 with per-cell clock-gating on zero-weight elements; and
demonstrate interoperability with existing ternary computing efforts via the
\texttt{ternlang-compat} bridge crate.

Beyond the execution substrate, we introduce two higher-level contributions.
The \textbf{Ternary AI Reasoning Toolkit} (Phase~8) provides five structural
primitives—iterative deliberation, coalition voting, multi-dimensional action
gating, scalar-to-temperature bridging, and hallucination scoring—that make
any AI agent's decision loop architecturally ternary rather than merely
prompted. The \textbf{MoE-13 Ternary Orchestrator} (Phase~9) implements a
Mixture-of-Experts head with dual-key synergistic routing, $1{+}1{=}3$
emergent triad synthesis, a hard-wired safety veto, and a three-tier memory
mesh, realising the MoE-13 architecture~\cite{moe13_2026} in 146 passing
tests across the full workspace. All nine crates are published to
\texttt{crates.io} under open-core licensing. A companion corpus of
twenty executable \texttt{.tern} programs demonstrates the hold state
across domains spanning aerospace, medicine, distributed systems, civic
governance, and AI agent design.
\end{abstract}

\begin{IEEEkeywords}
balanced ternary, ternary logic, sparse inference, BitNet, domain-specific
language, virtual machine, actor model, FPGA synthesis, Verilog,
mixture-of-experts, AI reasoning, ternary orchestration
\end{IEEEkeywords}

% ─────────────────────────────────────────────────────────────────────────────
\section{Introduction}
\label{sec:intro}

The computational substrate underlying modern artificial intelligence is
binary. Floating-point arithmetic, two-state memory, and Boolean logic gates
have driven five decades of progress—but they introduce a fundamental
representational mismatch when modelling systems that are inherently
three-valued: \emph{affirmed}, \emph{denied}, and \emph{undecided}.

Clinical diagnosis, legal reasoning, sensor fusion under noise, and
multi-agent consensus all require a native neutral state that binary computing
forces to encode as a special case—null pointers, NaN, sentinel values, or
probabilistic scores collapsed to a threshold. Each encoding is a workaround
for an absent primitive.

Balanced ternary~\cite{knuth1997art} provides that primitive. A \emph{trit}
$t \in \{-1, 0, +1\}$ carries three symmetric values. The neutral value $0$
is \emph{active}—a deliberate state of hold, not an empty bit pattern.
Balanced ternary arithmetic is self-complementing: negation requires no
special-case handling. And at the scale of modern neural networks, where
BitNet~\cite{ma2024era} and related work show that ternary-quantized weights
preserve accuracy with dramatically reduced computation, the case for a
ternary-native execution substrate is both theoretical and empirical.

Despite this, the ternary computing field remains fragmented: hobbyist
emulators, academic EDA tools for memristor hardware~\cite{bos2020ternary},
isolated Lisp interpreters, and hardware simulators without compiler support.
No project provides the full vertical stack.

\textbf{Ternlang} fills this gap. Our contributions are:

\begin{enumerate}
  \item A \emph{language design} for balanced ternary: three-way exhaustive
        pattern matching, first-class trit tensors, and an actor model for
        ternary message passing (\S\ref{sec:language}).
  \item The \textbf{BET ISA}: a formal 2-bit-encoded instruction set with
        opcode coverage from arithmetic through tensor operations and
        distributed agent control (\S\ref{sec:isa}).
  \item \textbf{TSPARSE\_MATMUL}: a VM opcode that skips zero-weight
        multiplications at the instruction level, realising the sparsity
        benefit of ternary quantization without software overhead
        (\S\ref{sec:sparse}).
  \item A \textbf{Verilog-2001 hardware backend} with synthesisable sparse
        matmul array and full BET processor, plus a cycle-accurate RTL
        simulator in pure Rust (\S\ref{sec:hdl}).
  \item A \textbf{Ternary AI Reasoning Toolkit}: five structural
        primitives—deliberation, coalition vote, action gate, scalar
        temperature, and hallucination score—that make any AI agent's
        decision loop architecturally ternary (\S\ref{sec:reasoning}).
  \item The \textbf{MoE-13 Ternary Orchestrator}: dual-key synergistic
        routing, $1{+}1{=}3$ emergent triad synthesis, a hard safety veto,
        and three-tier memory mesh, exposed as MCP tool calls
        (\S\ref{sec:moe}).
  \item \textbf{Ecosystem bridges} connecting existing ternary projects to
        the BET VM as a common runtime (\S\ref{sec:ecosystem}).
  \item A \textbf{reference example corpus}: twenty executable
        \texttt{.tern} programs demonstrating the hold state in aerospace,
        medicine, distributed systems, hardware pipelines, civic governance,
        finance, and AI agent design—published open-source alongside all
        nine workspace crates on \texttt{crates.io}.
\end{enumerate}

% ─────────────────────────────────────────────────────────────────────────────
\section{Background: Balanced Ternary}
\label{sec:background}

\subsection{Trit arithmetic}

A \emph{trit} $t \in T = \{-1, 0, +1\}$ participates in the following
complete operations~\cite{knuth1997art}:

\begin{definition}[Ternary addition with carry]
For trits $a, b \in T$, define $s = \mathrm{sum}(a,b)$ and
$c = \mathrm{carry}(a,b)$ such that $a + b = 3c + s$, where
$s, c \in T$.
\end{definition}

The complete sum/carry table is:

\begin{table}[h]
\centering
\small
\caption{Balanced ternary addition (sum, carry)}
\label{tab:add}
\begin{tabular}{c|ccc}
\toprule
$a \backslash b$ & $-1$ & $0$ & $+1$ \\
\midrule
$-1$ & $(+1,-1)$ & $(-1,\ 0)$ & $(0,\ 0)$ \\
$\ 0$ & $(-1,\ 0)$ & $(0,\ 0)$ & $(+1,\ 0)$ \\
$+1$ & $(0,\ 0)$ & $(+1,\ 0)$ & $(-1,+1)$ \\
\bottomrule
\end{tabular}
\end{table}

\begin{definition}[Ternary multiplication]
$\mathrm{mul}(a,b) = a \cdot b$ under integer multiplication,
restricted to $T$: $\mathrm{mul}(a,b) \in T$ since $|a|,|b| \leq 1$.
\end{definition}

\begin{definition}[Consensus (ternary OR)]
$\mathrm{cons}(a,b) = a$ if $a = b$, else $0$.
\end{definition}

\begin{definition}[Negation]
$\mathrm{neg}(t) = -t$. Note $\mathrm{neg}(0) = 0$.
\end{definition}

\subsection{The 2-bit BET encoding}

Hardware naturally operates in binary. We encode trits as 2-bit pairs:

\begin{table}[h]
\centering
\small
\caption{BET 2-bit trit encoding}
\label{tab:encoding}
\begin{tabular}{ccc}
\toprule
Bits & Value & Meaning \\
\midrule
\texttt{0b01} & $-1$ & conflict \\
\texttt{0b10} & $+1$ & truth \\
\texttt{0b11} & $\ 0$ & hold \\
\texttt{0b00} & --- & FAULT (invalid) \\
\bottomrule
\end{tabular}
\end{table}

This encoding is \emph{not} two's complement. The bit pattern \texttt{0b11}
for zero is chosen so that the all-ones reset state of a register file
initialises every register to hold—the semantically correct neutral value—
without special reset logic. Negation maps to a simple bit swap:
\texttt{0b01} $\leftrightarrow$ \texttt{0b10}, leaving \texttt{0b11}
invariant.

\begin{proposition}[Negation as bit swap]
For encoding $e(t)$ as above, $e(\mathrm{neg}(t)) = \{e(t)[0], e(t)[1]\}$
(swap bits).
\end{proposition}

% ─────────────────────────────────────────────────────────────────────────────
\section{The Ternlang Language}
\label{sec:language}

\subsection{Design principles}

Ternlang is statically typed, expression-oriented, and compiled to BET
bytecode. Three design decisions distinguish it from binary-native languages
adapted to ternary:

\textbf{Exhaustive three-way matching.} Every \texttt{match} expression must
cover all three trit arms. The compiler rejects non-exhaustive matches at
parse time, eliminating an entire class of runtime error present in languages
that bolt ternary logic onto binary pattern matching.

\textbf{$0$ is active.} The type system and semantics assign distinct meaning
to $-1$ (conflict), $0$ (hold, active neutral), and $+1$ (truth). There is no
null, no undefined, no NaN. A trit always carries a definite value.

\textbf{Sparsity is a language feature.} The \texttt{@sparseskip} directive
marks a tensor operation as sparse-aware, routing the compiler to emit
\texttt{TSPARSE\_MATMUL} instead of \texttt{TMATMUL}. Sparsity is expressed
in the source language, not discovered by the optimizer.

\subsection{Core language constructs}

\begin{lstlisting}[caption={Ternary classifier with exhaustive match}]
fn classify(signal: trit) -> trit {
    match signal {
        -1 => conflict()  // disagreement confirmed
         0 => hold()      // awaiting evidence
        +1 => truth()     // assertion confirmed
    }
}
\end{lstlisting}

\begin{lstlisting}[caption={Sparse matrix multiply with @sparseskip}]
// Route to TSPARSE_MATMUL at the ISA level
@sparseskip
let output: trittensor<8 x 8> =
    matmul(input, weights);
\end{lstlisting}

\begin{lstlisting}[caption={Actor model for ternary message passing}]
agent Voter {
    fn handle(msg: trit) -> trit {
        consensus(msg, hold())
    }
}

let v: agentref = spawn Voter;
send v truth();
let decision: trit = await v;
\end{lstlisting}

\subsection{Type system}

The core types are \texttt{trit} (a single balanced ternary value),
\texttt{trittensor<N x M>} (an $N \times M$ matrix of trits stored on the
tensor heap), and \texttt{agentref} (a handle to a running actor instance,
local or remote). Struct types with trit/tensor fields are supported via
field-name mangling in the register allocator.

\subsection{Ternary Error Propagation}
\label{subsec:propagate}

Ternlang introduces a postfix \texttt{?} operator for ternary-native error
propagation, analogous to Rust's \texttt{?} operator but grounded in the
three-valued semantics of the trit space. For any expression $e$ of type
\texttt{trit}, writing $e\texttt{?}$ in a function body has the following
semantics:

\begin{equation}
  e\texttt{?} \;\triangleq\;
  \begin{cases}
    \text{return}\;{-1} & \text{if } e = {-1} \\
    e                   & \text{otherwise}
  \end{cases}
  \label{eq:propagate}
\end{equation}

This elevates trit $-1$ (conflict) from a value to a structural signal: any
function that receives a conflict can propagate it up the call chain without
explicit match arms. At the BET level, \texttt{expr?} compiles to:

\begin{lstlisting}[caption={BET bytecode for \texttt{expr?}}]
; expr already on stack: [val]
TDUP          ; [val, dup]
TJMP_NEG L1   ; consume dup; if -1 -> L1; else [val] continues
TJMP L2       ; skip early return
L1: TRET      ; return -1 to caller
L2: (continue with val on stack)
\end{lstlisting}

The parser disambiguates \texttt{?} from the ternary branching operator
(also \texttt{?}) using two-token lookahead: \texttt{expr?} followed by
\texttt{\{} is an uncertain branch; \texttt{expr?} followed by any other
token is a propagation operator.

\begin{lstlisting}[caption={Error propagation in a validation chain}]
fn validate_signal(x: trit) -> trit {
    let checked: trit = boundary_check(x)?; // propagate if conflict
    return range_check(checked)?;           // propagate again
}

fn main() -> trit {
    return validate_signal(1)?; // -1 if any check fails
}
\end{lstlisting}

\subsection{Cross-File Module System}
\label{subsec:modules}

Ternlang's module system supports two resolution strategies, tried in order:

\begin{enumerate}
  \item \textbf{Built-in stdlib} (compile-time embedded): \texttt{std::trit},
        \texttt{std::math}, \texttt{std::tensor}, \texttt{std::io},
        \texttt{ml::quantize}, \texttt{ml::inference} are embedded in the
        compiler binary via \texttt{include\_str!}—zero filesystem I/O at
        runtime.
  \item \textbf{User-defined modules}: a \texttt{ModuleResolver} walks
        \texttt{use} paths relative to the source file's directory. A
        declaration \texttt{use agents::voter;} resolves to
        \texttt{<source\_dir>/agents/voter.tern}, parsed and merged into
        the program before semantic analysis.
\end{enumerate}

Both strategies deduplicate: a module loaded by multiple \texttt{use}
statements is parsed exactly once, and functions already present in the
program are never duplicated. This makes \texttt{StdlibLoader::resolve()}
idempotent and safe to call multiple times.

\begin{lstlisting}[caption={Cross-file module import}]
// agents/voter.tern
fn weighted_vote(a: trit, b: trit, c: trit) -> trit {
    consensus(consensus(a, b), c)
}

// main.tern
fn main() -> trit {
    use agents::voter;
    return weighted_vote(1, 0, -1);
}
\end{lstlisting}

\subsection{Diagnostic Philosophy}
\label{subsec:diagnostics}

Ternlang's error messages carry structured codes and ternary-philosophical
commentary. Every diagnostic has a machine-readable code (for tooling and
documentation lookup) and a human-readable nudge that frames the error in
the language's conceptual model:

\begin{itemize}
  \item \texttt{[PARSE-001]} Unexpected token — \emph{"the lexer hit something
        it didn't expect. Check your syntax."}
  \item \texttt{[PARSE-004]} Non-exhaustive match — \emph{"ternary has three
        states — cover all three or the compiler won't let you through."}
  \item \texttt{[SCOPE-001]} Undefined variable — \emph{"hold state — declare
        before use."}
  \item \texttt{[FN-002]} Return type mismatch — \emph{"ternary contracts are
        strict."}
  \item \texttt{[BET-001]} Stack underflow — \emph{"you tried to pop a truth
        that wasn't there."}
  \item \texttt{[PROP-001]} \texttt{?} on non-trit — \emph{"only trit-returning
        functions can signal conflict. The third state requires a trit."}
\end{itemize}

This design ensures that newcomers encountering their first ternary error
receive both actionable guidance and an introduction to the philosophy
underpinning the type system.

% ─────────────────────────────────────────────────────────────────────────────
\section{The BET Instruction Set Architecture}
\label{sec:isa}

\subsection{Machine model}

The BET VM is a stack-based machine with:

\begin{itemize}
  \item \textbf{27 registers} (2 bits each), reset to \texttt{0b11} (hold)
  \item \textbf{Value stack}: unbounded, stores \texttt{Value} tagged union
  \item \textbf{Tensor heap}: indexed array of $N \times M$ trit matrices
  \item \textbf{Call stack}: return address stack for \texttt{TCALL}/\texttt{TRET}
  \item \textbf{Agent table}: maps type IDs to handler addresses and mailboxes
  \item \textbf{Carry register}: overflow from \texttt{TADD}
\end{itemize}

The choice of 27 registers reflects the ternary motif: $27 = 3^3$.

\subsection{Instruction encoding}

Instructions are variable-length, with a 1-byte opcode followed by 0–2
operand bytes. Table~\ref{tab:isa} lists the complete opcode map.

\begin{table}[h]
\centering
\small
\caption{BET ISA opcode reference (selected)}
\label{tab:isa}
\begin{tabular}{lll}
\toprule
Opcode & Mnemonic & Stack effect \\
\midrule
\texttt{0x00} & THALT            & --- \\
\texttt{0x01} & TPUSH \textit{t} & $\to t$ \\
\texttt{0x02} & TADD             & $a\ b \to s\ c$ \\
\texttt{0x03} & TMUL             & $a\ b \to t$ \\
\texttt{0x04} & TNEG             & $t \to \mathrm{neg}(t)$ \\
\texttt{0x05} & TJMP\_POS \textit{addr} & $t \to$ (jmp if $t=+1$) \\
\texttt{0x06} & TJMP\_ZERO \textit{addr}& $t \to$ (jmp if $t=0$) \\
\texttt{0x07} & TJMP\_NEG \textit{addr} & $t \to$ (jmp if $t=-1$) \\
\texttt{0x08} & TSTORE \textit{r}& $t \to$ (reg[$r$]$\leftarrow t$) \\
\texttt{0x09} & TLOAD \textit{r} & $\to$ reg[$r$] \\
\texttt{0x0b} & TJMP \textit{addr}& --- \\
\texttt{0x0e} & TCONS            & $a\ b \to \mathrm{cons}(a,b)$ \\
\texttt{0x0f} & TALLOC \textit{N M} & $\to$ ref \\
\texttt{0x10} & TCALL \textit{addr} & (push PC; jmp) \\
\texttt{0x11} & TRET            & (pop PC; jmp) \\
\texttt{0x20} & TMATMUL         & $r_A\ r_B \to r_C$ \\
\texttt{0x21} & TSPARSE\_MATMUL & $r_A\ r_B \to r_C$ \\
\texttt{0x22} & TIDX            & ref $i\ j \to t$ \\
\texttt{0x23} & TSET            & ref $i\ j\ t \to$ \\
\texttt{0x24} & TSHAPE          & ref $\to N\ M$ \\
\texttt{0x25} & TSPARSITY       & ref $\to$ count \\
\texttt{0x30} & TSPAWN \textit{id} & $\to$ agentref \\
\texttt{0x31} & TSEND           & agentref $m \to$ \\
\texttt{0x32} & TAWAIT          & agentref $\to t$ \\
\bottomrule
\end{tabular}
\end{table}

% ─────────────────────────────────────────────────────────────────────────────
\section{Sparse Ternary Inference}
\label{sec:sparse}

\subsection{Ternary quantization}

BitNet-style ternary weight quantization~\cite{ma2024era} maps floating-point
weights $w \in \mathbb{R}$ to $\hat{w} \in \{-1, 0, +1\}$ using a threshold
$\tau$:

\begin{equation}
  \hat{w} = \begin{cases}
    +1 & w >  \tau \\
     0 & |w| \leq \tau \\
    -1 & w < -\tau
  \end{cases}
  \qquad \tau = \frac{1}{2} \cdot \mathbb{E}[|w|]
  \label{eq:quantize}
\end{equation}

The resulting weight distribution is heavily concentrated at $0$ (hold):
typical language model weights at BitNet scale show $55$–$65\%$ zero elements
after quantization with $\tau$ as defined in~(\ref{eq:quantize}).

\subsection{The @sparseskip directive and TSPARSE\_MATMUL}

The key insight is that a multiply against a zero-weight trit produces zero
regardless of the input value:

\begin{equation}
  \mathrm{mul}(a, 0) = 0 \quad \forall a \in T
  \label{eq:zero-mul}
\end{equation}

In a dense $N \times M$ matrix multiply, every element contributes
$N \cdot M$ multiplications. After ternary quantization with sparsity $\rho$
(fraction of zero-weight elements), only $(1-\rho) \cdot N \cdot M$
multiplications have non-trivial results.

\texttt{TSPARSE\_MATMUL} (opcode \texttt{0x21}) implements a sparse
inner-product loop that tests the weight trit before executing the multiply:

\begin{lstlisting}[language=C, caption={TSPARSE\_MATMUL pseudocode}]
for i in 0..N:
  for j in 0..M:
    for k in 0..K:
      w = W[k][j]
      if w == HOLD: continue   // skip
      acc[i][j] += mul(A[i][k], w)
\end{lstlisting}

The directive \texttt{@sparseskip} in the source language is a compile-time
annotation that causes the codegen to emit \texttt{TSPARSE\_MATMUL} instead of
\texttt{TMATMUL} for the following \texttt{matmul()} call. Sparsity awareness
is thus a \emph{source-language} property surfaced at the ISA level, not a
runtime optimization that may or may not trigger.

\subsection{Benchmark results}

We evaluate sparse versus dense matrix multiply on $512 \times 512$ weight
matrices quantized from normally distributed $\mathcal{N}(0,1)$ float weights
using threshold~(\ref{eq:quantize}).

\begin{table}[h]
\centering
\small
\caption{Sparse vs.\ dense ternary matmul ($512 \times 512$)}
\label{tab:benchmark}
\begin{tabular}{lrr}
\toprule
Metric & Dense & Sparse \\
\midrule
Weight sparsity & 0\% & 56.2\% \\
Multiply operations & 262,144 & 115,343 \\
Skipped operations & 0 & 146,801 \\
Speedup & $1.0\times$ & $\mathbf{2.27\times}$ \\
\bottomrule
\end{tabular}
\end{table}

The $2.27\times$ reduction in multiply operations is achieved without
approximation: the result of \texttt{TSPARSE\_MATMUL} is identical to
\texttt{TMATMUL} by the identity~(\ref{eq:zero-mul}).

% ─────────────────────────────────────────────────────────────────────────────
\section{Hardware Backend}
\label{sec:hdl}

\subsection{Verilog-2001 primitives}

The \texttt{ternlang-hdl} crate generates synthesisable Verilog-2001 modules
from the BET ISA. Each trit becomes a \texttt{[1:0]} bus. The core primitives
are:

\begin{itemize}
  \item \texttt{trit\_neg}: bit swap (\texttt{assign y = \{a[0], a[1]\}})
  \item \texttt{trit\_cons}: \texttt{assign y = (a == b) ? a : 2'b11}
  \item \texttt{trit\_mul}: zero-skip detect with combinational multiply
  \item \texttt{trit\_add}: 9-case Verilog \texttt{case} with carry
  \item \texttt{trit\_reg}: synchronous D-register with reset-to-hold
  \item \texttt{bet\_alu}: top-level ALU selecting among primitives by op code
\end{itemize}

\subsection{Sparse matmul array}

The synthesisable sparse matmul array instantiates an $N \times N$ grid of
processing elements. Each cell contains a weight register and a clock-gate
signal:

\begin{lstlisting}[language=Verilog,
  caption={Per-cell clock gating in sparse matmul}]
wire [1:0] w_ij = weight_reg[i][j];
wire skip = (w_ij == 2'b11);  // hold = zero weight
wire [1:0] contrib = skip ?
    2'b11 :                    // hold if skipped
    trit_mul(a_i, w_ij);
\end{lstlisting}

Clock-gating on the \texttt{skip} signal prevents switching activity in
zero-weight cells, delivering dynamic power reduction proportional to sparsity.

\subsection{BET processor}

The full \texttt{bet\_processor} module wires together:
\begin{itemize}
  \item \texttt{bet\_regfile}: 27-register file ($27 \times 2$ bits)
  \item \texttt{bet\_pc}: 16-bit program counter with load port for jumps
  \item \texttt{bet\_control}: single-cycle decode, all 51 opcodes mapped to
        control signals (alu\_op, reg\_we, mem\_re, mem\_we, pc\_load,
        is\_call, is\_ret, is\_halt, is\_spawn, is\_send, is\_await)
\end{itemize}

An Icarus Verilog simulation wrapper (\texttt{BetSimEmitter}) generates
complete self-contained testbenches from BET bytecode, including inline ROM
initialization, reset sequencing, and VCD waveform export for GTKWave.

% ─────────────────────────────────────────────────────────────────────────────
\section{Native Compilation: AST-to-C Backend}
\label{sec:codegen}

The \texttt{ternlang-codegen} crate provides a second compilation target
alongside BET bytecode: a self-contained C11 source file that can be compiled
with any standard C compiler. This opens three new deployment paths:
native-speed execution, cross-compilation to embedded targets, and human-readable
output for inspection and security audit.

\subsection{Transpilation Architecture}

The \texttt{CTranspiler} operates directly on the abstract syntax tree (Program)
produced by the parser, bypassing the BET VM entirely:

\begin{equation}
  \texttt{.tern} \xrightarrow{\text{parse}} \text{AST}
  \xrightarrow{\text{CTranspiler}} \texttt{.c}
  \xrightarrow{\text{gcc/clang}} \text{native binary}
  \label{eq:ctranspile}
\end{equation}

The generated C file is fully self-contained: a header of inline trit
primitives using \texttt{int8\_t} with values $\{-1, 0, +1\}$ is prepended,
covering \texttt{trit\_add}, \texttt{trit\_mul}, \texttt{trit\_neg},
\texttt{trit\_consensus}, \texttt{trit\_abs}, and the built-in constants
\texttt{trit\_truth()}, \texttt{trit\_hold()}, \texttt{trit\_conflict()}.

\subsection{Language Construct Mapping}

The transpiler maps every ternlang construct to its natural C equivalent:

\begin{itemize}
  \item \texttt{match} $\rightarrow$ \texttt{switch (int)} with \texttt{case -1},
        \texttt{case 0}, \texttt{case 1} arms
  \item \texttt{struct} definitions $\rightarrow$ \texttt{typedef struct \{ \ldots \}}
  \item \texttt{fn} $\rightarrow$ C function with forward declaration (enables
        mutual recursion without re-ordering)
  \item \texttt{loop / while / for} $\rightarrow$ \texttt{for(;;)},
        \texttt{while(1)} with inner ternary condition dispatch
  \item \texttt{expr?} $\rightarrow$ \texttt{\_\_TERN\_PROPAGATE(expr)} macro
        slot, enabling the caller to define the early-return idiom in C
  \item \texttt{main()} $\rightarrow$ \texttt{tern\_main()} (to avoid clashing
        with C's entry point); a generated \texttt{main()} calls \texttt{tern\_main()}
        and translates the trit result to an exit code
\end{itemize}

\begin{lstlisting}[caption={Ternlang source and generated C output}]
// ternlang source
fn decide(x: trit) -> trit {
    match x {
        -1 => conflict()
         0 => hold()
        +1 => truth()
    }
}

// generated C
trit decide(trit x) {
    switch ((int)x) {
        case -1: return trit_conflict(); break;
        case  0: return trit_hold();     break;
        case  1: return trit_truth();    break;
    }
}
\end{lstlisting}

The C backend does not yet support the actor model (\texttt{spawn/send/await})
or tensor heap operations, which remain BET-VM specific. These emit comments
in the output, making the generated file auditable even for unsupported constructs.

% ─────────────────────────────────────────────────────────────────────────────
\section{Actor Model and Distributed Runtime}
\label{sec:actors}

\subsection{Local actors}

Ternlang implements the actor model via three ISA primitives:
\texttt{TSPAWN} creates an agent instance and returns an \texttt{agentref};
\texttt{TSEND} enqueues a trit message in the agent's mailbox; \texttt{TAWAIT}
dequeues a message, invokes the handler, and returns the trit result. The
mailbox is a FIFO (\texttt{VecDeque<Value>}) allocated per agent instance.

\subsection{Distributed actors}

The \texttt{ternlang-runtime} crate extends the actor model to multi-node
systems via TCP transport. A \texttt{TernNode} binds a TCP port, maintains a
peer connection map, and exposes \texttt{remote\_send}/\texttt{remote\_await}
over a newline-delimited JSON wire protocol:

\begin{lstlisting}[language=ternlang,
  caption={Remote actor spawn in ternlang}]
let remote_voter: agentref =
    spawn remote "192.168.1.42:7373" Voter;
send remote_voter truth();
let r: trit = await remote_voter;
\end{lstlisting}

The wire protocol uses four message types (\texttt{Send}, \texttt{Await},
\texttt{Reply}, \texttt{Error}) serialised as tagged JSON objects, enabling
interoperability with non-Rust implementations.

% ─────────────────────────────────────────────────────────────────────────────
\section{Ternary AI Reasoning Toolkit}
\label{sec:reasoning}

Modern AI agents are structurally binary in their decision loops: evidence
is collapsed to a scalar, thresholded, and converted to a Boolean act/wait
flag. The three-valued nature of evidence—affirm, hold, reject—is treated
as a side-effect of confidence scores rather than a first-class type.

Phase~8 of the TIS introduces five primitives in \texttt{ternlang-ml} that
make the decision loop architecturally ternary.

\subsection{Deliberation Engine}

The \texttt{DeliberationEngine} models iterative evidence accumulation using
an exponential moving average (EMA):

\begin{equation}
  S_r = \alpha \cdot e_r + (1-\alpha) \cdot S_{r-1},
  \quad \alpha \in (0, 1]
  \label{eq:ema}
\end{equation}

where $e_r$ is the mean of new evidence signals in round~$r$. The engine
commits to a ternary decision only when the confidence of $S_r$ exceeds a
target threshold; otherwise it remains in the hold state and requests a
further evidence round. This models the human ``let me think about it''
behaviour—a native ternary computation rather than a binary timeout.

\subsection{Coalition Vote}

\texttt{coalition\_vote()} aggregates $N$ independent agent verdicts—each a
trit with an associated confidence and weight—into a single result with
quorum, dissent, and abstain statistics:

\begin{equation}
  \hat{t} = \mathrm{sign}\!\left(
    \frac{\sum_i w_i \cdot c_i \cdot t_i}{\sum_i w_i \cdot c_i}
  \right)
  \label{eq:coalition}
\end{equation}

where $t_i \in \{-1,0,+1\}$, $c_i \in [0,1]$ is confidence, and
$w_i > 0$ is importance weight.

\subsection{Action Gate}

The action gate enforces a structural separation between hard constraints
and soft preferences. Each \texttt{GateDimension} carries an evidence
signal, a weight, and an optional \texttt{hard\_block} flag. A hard-block
dimension with negative evidence vetoes the action unconditionally,
regardless of all other dimensions:

\begin{equation}
  \text{verdict} = \begin{cases}
    \text{Block} & \exists\, d_i : \text{hard\_block}(d_i) \wedge t_{d_i} = -1 \\
    \text{sign}(\bar{S}) & \text{otherwise}
  \end{cases}
  \label{eq:gate}
\end{equation}

This is the structural analogue of a safety interlock: the veto is
unconditional, not probabilistic.

\subsection{Scalar Temperature and Hallucination Score}

\texttt{scalar\_temperature()} bridges the ternary decision to the LLM
sampling temperature domain: affirm at high confidence maps to a low
temperature (focused generation), hold maps to a mid temperature
(exploratory generation), and reject maps to near-zero (cautious refusal).

\texttt{hallucination\_score()} maps the variance of an evidence signal
vector to a trust trit: low variance (consistent signals) yields trust $+1$;
high variance yields $-1$, signalling that the agent's signals are
incoherent and should not be acted upon.

% ─────────────────────────────────────────────────────────────────────────────
\section{MoE-13 Ternary Orchestrator}
\label{sec:moe}

The MoE-13 architecture~\cite{moe13_2026} is implemented in the
\texttt{ternlang-moe} crate. It realises a ternary Mixture-of-Experts head
whose routing, synthesis, safety gate, and memory are all expressed natively
in balanced ternary semantics.

\subsection{Six-Dimensional Competence Space}

Each expert is characterised by a \emph{competence vector}
$\mathbf{v} \in [-1,1]^6$ across six axes:

\begin{equation}
  \mathbf{v} = \bigl[
    v_\text{syntax},\,
    v_\text{world},\,
    v_\text{reason},\,
    v_\text{tool},\,
    v_\text{persona},\,
    v_\text{safety}
  \bigr]
  \label{eq:competence}
\end{equation}

The synergy between two experts is defined as the complement of their cosine
similarity—orthogonal experts are maximally complementary:

\begin{equation}
  \sigma(\mathbf{v}_i, \mathbf{v}_j)
    = \frac{1 - \cos(\mathbf{v}_i, \mathbf{v}_j)}{2}
    \;\in\; [0, 1]
  \label{eq:synergy}
\end{equation}

\subsection{Dual-Key Synergistic Routing}

For each candidate expert pair $(i, j)$, the routing score combines
relevance to the query vector $\mathbf{q}$ with inter-expert synergy:

\begin{equation}
  \mathrm{score}(i, j)
    = \rho_i(\mathbf{q}) \cdot \rho_j(\mathbf{q}) \cdot \sigma(\mathbf{v}_i, \mathbf{v}_j)
  \label{eq:routing}
\end{equation}

where $\rho_k(\mathbf{q}) = \max(0,\, \cos(\mathbf{v}_k, \mathbf{q}))$
is the relevance of expert~$k$ to query $\mathbf{q}$. The pair with the
highest score is selected. Crucially, two highly relevant but similar
experts are penalised by low $\sigma$; the router favours complementary
coverage over redundant strength.

\subsection{$1{+}1{=}3$ Triad Synthesis}

After expert evaluation, an emergent \emph{triad field} is synthesised from
the two selected competence vectors and the routing synergy:

\begin{equation}
  \mathbf{E}_k = \sigma_{ij} \cdot \frac{\mathbf{v}_i + \mathbf{v}_j}{2}
  \label{eq:triad}
\end{equation}

This is the formal expression of $1{+}1{=}3$: two expert signals at high
synergy produce a third signal—the emergent field $\mathbf{E}_k$—that
neither expert could produce in isolation. The field modulates the
subsequent vote and the LLM temperature hint.

\subsection{Safety Hard Gate and Axis-6 Veto}

Dimension~6 of every competence vector is the \emph{safety axis}. Before
any vote is computed, the safety projection of the triad field is evaluated:

\begin{equation}
  \text{veto} \iff E_{k,\text{safety}} < \theta_s
  \label{eq:veto}
\end{equation}

where $\theta_s$ is the configurable safety threshold (default $-0.3$).
A veto immediately returns trit $-1$ at confidence $1.0$ and logs the event
to the Axis-tier memory for audit. No downstream reasoning can override it.

\subsection{Hold State and Tiebreaker}

When the weighted vote yields trit $0$ or aggregate confidence falls below
a hold threshold, the orchestrator does not immediately return. It invokes a
tiebreaker expert—selected by highest reasoning-axis score among inactive
experts—and re-votes with up to $\max\_\text{active}=4$ experts total. Only
if the tiebreaker also fails to resolve does the orchestrator return the
hold state to the caller, signalling that further evidence is needed.

\subsection{Three-Tier Memory Mesh}

The orchestrator maintains three memory tiers:

\begin{itemize}
  \item \textbf{Node} (TTL: seconds): LRU-bounded volatile store for
        within-session context. Evicts on capacity overflow.
  \item \textbf{Cluster} (TTL: minutes): routing frequency counters enabling
        mode-collapse detection ($\text{risk} = \max\_\text{pair} / \text{total}$).
  \item \textbf{Axis} (persistent): immutable audit log of safety vetoes with
        timestamp, expert identity, and query hash. Global priors over expert
        performance.
\end{itemize}

\subsection{Nine-Step Inference Pipeline}

The full orchestration pass executes in nine deterministic steps:
(1)~encode query to evidence vector;
(2)~dual-key route to best expert pair;
(3)~evaluate both experts independently;
(4)~synthesise triad field via~(\ref{eq:triad});
(5)~safety hard gate via~(\ref{eq:veto});
(6)~weighted trit vote with synergy amplification;
(7)~hold detection;
(8)~optional tiebreaker invocation;
(9)~return \texttt{OrchestrationResult} and update all three memory tiers.

\subsection{Thirteen-Agent Dual-Signal Deliberation}
\label{subsec:harness}

The \texttt{AgentHarness} in \texttt{ternlang-moe/src/agents/} extends the
9-step orchestration pass with a structured deliberation layer of 13
specialised agents, each computing two independent signals:

\begin{equation}
  t_i = \begin{cases}
    +1 & \text{if positive signal} \geq \theta^+ \text{ and negative signal} < \theta^- \\
    -1 & \text{if negative signal} \geq \theta^- \text{ and positive signal} < \theta^+  \\
    \phantom{+}0 & \text{(stasis: signals in equilibrium)}
  \end{cases}
  \label{eq:dual_signal}
\end{equation}

where $\theta^+$ and $\theta^-$ are agent-specific thresholds. The 13 agents cover:
\textbf{Syntax} (structural token analysis, bracket balance),
\textbf{WorldKnowledge} (domain vocabulary density, proper noun frequency),
\textbf{DeductiveReason} (premise + conclusion marker co-occurrence),
\textbf{InductiveReason} (example density, generalisation leap detection),
\textbf{ToolUse} (imperative verb strength, passive construction penalty),
\textbf{Persona} (depersonalisation detection),
\textbf{Safety} (hard risk keywords $\to$ trit $-1$ at confidence 0.99),
\textbf{FactCheck} (overconfidence markers $\to$ hallucination flag),
\textbf{CausalReason} (requires both cause and effect markers),
\textbf{AmbiguityRes} (hard-vague $\geq 2$ markers $\to$ trit $-1$),
\textbf{MathReason} (impossible operations e.g.\ divide-by-zero $\to$ trit $-1$),
\textbf{ContextMem} (fresh-start signals vs.\ anaphoric reference density),
and \textbf{MetaSafety} (injection pattern detection at confidence 0.99).

The dual-signal design is philosophically significant: trit $0$ is not a
default or a fallback from a failed comparison. It is \emph{active stasis}—the
system's honest declaration that positive and negative evidence are balanced
and further information is required before committing.

\subsection{Introspective Hold: The Stable Attractor}
\label{subsec:hold}

The \texttt{run\_introspective()} method on \texttt{AgentHarness} implements a
\emph{stable attractor} for the hold state. Once reached, a stable hold is
permanent within the deliberation turn—it cannot be overridden by additional
tiebreaker invocations. The stable attractor condition is:

\begin{equation}
  \text{stable\_hold} \iff
  \lvert \text{affirm\_count} - \text{conflict\_count} \rvert \leq 1
  \;\wedge\;
  \text{engaged} \geq 4
  \label{eq:stable_hold}
\end{equation}

This formalises the philosophical claim: when at least four independent
agents have deliberated and the signal balance between affirmation and
conflict is indistinguishable, the appropriate epistemic response is to
hold—not to choose arbitrarily between $+1$ and $-1$. The hold state
is returned with \texttt{is\_stable\_hold: true} and a human-readable
\texttt{hold\_reason} string.

The aggregate verdict structure carries full deliberation metadata:

\begin{lstlisting}[caption={AggregateVerdict fields}]
pub struct AggregateVerdict {
    pub trit:           i8,
    pub confidence:     f32,
    pub verdicts:       Vec<ExpertVerdict>,
    pub is_stable_hold: bool,
    pub hold_reason:    Option<String>,
    pub affirm_count:   usize,
    pub conflict_count: usize,
    pub hold_count:     usize,
}
\end{lstlisting}

Safety and meta-safety form a hard gate: the evidence vector produced by
\texttt{to\_evidence\_vector()} maps axis~6 to
$\min(\text{Safety.confidence},\, \text{MetaSafety.confidence})$, ensuring
the safety dimension is always the most conservative signal in the system.

\subsection{Orchestrate-Full: Thirteen-Substage Pipeline}
\label{subsec:orchestrate_full}

The \texttt{orchestrate\_full()} method on \texttt{TernMoeOrchestrator}
composes the 13-agent deliberation with the 9-step routing pipeline into a
unified end-to-end pass:

\begin{enumerate}
  \item Run all 13 agents via \texttt{AgentHarness::run\_introspective()}
  \item Check for stable attractor hold — return immediately if true
  \item Map 13 verdicts to 6D evidence vector via \texttt{to\_evidence\_vector()}
  \item Check safety hard gate on safety axis before MoE routing
  \item Run standard 9-step MoE orchestration with enriched evidence
  \item Return \texttt{OrchestrationResult} with stable hold flag propagated
\end{enumerate}

This creates a two-tier system: the fast ternary language models (13 agents,
keyword-level deliberation) form a pre-filter that enriches the evidence
before the slower but more semantically powerful MoE routing pass. Most
queries that are clearly safe and unambiguous complete in the agent tier;
only complex or borderline queries reach the full routing pipeline.

\subsection{MCP Exposure}

The orchestrator is exposed as three MCP tool calls:
\texttt{moe\_orchestrate} runs a full pass and returns the complete result
including routing pair, triad field, verdicts, temperature, and prompt hint;
\texttt{moe\_deliberate} wraps the EMA deliberation engine for multi-round
iterative reasoning; \texttt{trit\_action\_gate} exposes the hard-block gate
as a standalone safety check. Any MCP-compatible AI client connected to
\texttt{ternlang-mcp} gains access to the full ternary reasoning stack as
tool calls.

% ─────────────────────────────────────────────────────────────────────────────
\section{Qutrit Neural Networks and Ternlang}
\label{sec:qnn}

The Qutrit Neural Network (QNN) framework~\cite{kepp2026qnn} establishes a
formal correspondence between quantum qutrit states and balanced ternary values.
The three qutrit basis states $\lvert -\rangle$ (decay), $\lvert 0\rangle$
(stasis), and $\lvert +\rangle$ (growth) map exactly onto the ternlang trit
space:

\begin{equation}
  \lvert -\rangle \leftrightarrow -1, \quad
  \lvert 0\rangle \leftrightarrow 0,  \quad
  \lvert +\rangle \leftrightarrow +1
  \label{eq:qnn_map}
\end{equation}

This correspondence is not coincidental. Balanced ternary's symmetric
structure around zero and its active neutral state mirror the physical
properties of three-state quantum systems, making ternlang a natural
implementation substrate for QNN inference.

\subsection{Tesseract Recursive Syntax Stabilisation}

The QNN paper introduces \emph{Tesseract Recursive Syntax} (TRS), a
stabilisation mechanism that prevents qutrit networks from collapsing
into binary modes under iterative refinement. In ternlang, TRS corresponds
to the introspective hold architecture (Section~\ref{subsec:hold}): when
affirmation and conflict signals are balanced, the system stabilises at
trit~$0$ rather than forcing a premature commitment.

\subsection{QNN Example Suite}

The ternlang repository contains 15 QNN-inspired \texttt{.tern} example
programs (files 251–265) covering:

\begin{itemize}
  \item Qutrit superposition encoding and collapse simulation
  \item Ternary attention mechanisms with $\{-1, 0, +1\}$ key-query products
  \item Qutrit activation functions replacing ReLU/GELU in ternary networks
  \item Phase-encoding of temporal signals as trit sequences
  \item TRS stabilisation loops using ternlang's \texttt{loop} construct
        and the \texttt{?} propagation operator for collapse detection
  \item Qutrit gradient approximation via balanced ternary finite differences
\end{itemize}

These examples demonstrate that ternlang's syntax is sufficient to express
QNN computations natively, without any binary adaptor layer. The \texttt{@sparseskip}
directive applies directly to qutrit weight matrices: at realistic qutrit
sparsity levels ($\approx 33$\% zero states at initialisation), the
\texttt{TSPARSE\_MATMUL} kernel provides the same $8$--$14\times$ speedup
over dense computation observed in classical BitNet workloads.

% ─────────────────────────────────────────────────────────────────────────────
\section{Related Work}
\label{sec:related}

\textbf{Balanced ternary computing.} Knuth~\cite{knuth1997art} provides the
mathematical foundation. The Setun computer (1958)~\cite{brousentsov2002setun}
demonstrated physical ternary hardware. Modern revivals include academic
hardware studies at the University of South-Eastern Norway~\cite{bos2020ternary},
which targets C-to-ternary compilation for EDA tools—an approach that inherits
binary-native semantics and misses the symmetric properties of balanced ternary
that ternlang is designed to exploit natively.

\textbf{Ternary emulators.} Several open-source projects implement ternary
VMs in Python or JavaScript (Brandon Smith's 9-trit RISC simulator; Owlet,
an S-expression ternary interpreter in Node.js). These provide single-layer
implementations without compiler, ML kernels, or hardware backends.
\texttt{ternlang-compat} provides assembler-level bridges to both.

\textbf{Ternary quantization for neural networks.} BitNet~\cite{ma2024era}
and subsequent work (e.g., BitNet b1.58~\cite{ma2024bitnet158}) demonstrate
that large language models can be trained with ternary weights while retaining
competitive perplexity. The present work is the first to surface this property
as a first-class ISA opcode (\texttt{TSPARSE\_MATMUL}) rather than a software
library optimization.

\textbf{Quantum ternary.} Qutrits (3-level quantum systems) map naturally to
trit values $\{|{-1}\rangle, |0\rangle, |{+1}\rangle\}$. The BET encoding and
\texttt{trittensor} type system are structurally compatible with qutrit state
spaces; we leave the formal mapping to future work.

% ─────────────────────────────────────────────────────────────────────────────
\section{The Ternary Computing Ecosystem}
\label{sec:ecosystem}

A stated goal of ternlang is to serve as a \emph{convergence point} for the
fragmented ternary computing field. Table~\ref{tab:ecosystem} maps active
projects to their interoperability path via \texttt{ternlang-compat}.

\begin{table}[h]
\centering
\small
\caption{Ecosystem compatibility map}
\label{tab:ecosystem}
\begin{tabular}{p{2.2cm}p{2.2cm}p{2.8cm}}
\toprule
Project & Interface & Bridge \\
\midrule
Brandon Smith 9-trit sim & \texttt{.tasm} assembly & \texttt{TasmAssembler} $\to$ BET bytecode \\
Owlet (S-expr) & S-expression syntax & \texttt{OwletParser} $\to$ ternlang AST \\
USN / EDA tools & C-to-ternary & Academic whitepaper; ISA interop \\
Trit-Rust crate & Rust i8 trits & Superseded by ternlang-core \\
Q-Ternary & Qutrits & \texttt{trittensor} state model (future) \\
\bottomrule
\end{tabular}
\end{table}

Beyond runtime interoperability, ternlang aims to serve as a
\emph{conceptual archive} for prior ternary computing work.
Brandon Smith's Python 9-trit RISC simulator demonstrated a complete
fetch-decode-execute cycle in balanced ternary—an existence proof that
received little downstream adoption due to the absence of a compiler
ecosystem. The Owlet interpreter proved that S-expression evaluation maps
cleanly to a three-valued substrate. Both contributions are honoured in
the \texttt{examples/} corpus: \texttt{09\_risc\_fetch\_decode.tern}
translates Smith's pipeline decision logic into native BET language syntax,
and \texttt{13\_owlet\_bridge.tern} models cooperative ternary evaluation
with hold-as-suspension semantics. In this way, prior work that was
``slowly forgotten'' acquires a working executable form—runnable on the
same VM that drives the MoE-13 orchestrator.

% ─────────────────────────────────────────────────────────────────────────────
\section{Implementation Status}
\label{sec:status}

Ternlang is implemented in Rust and comprises the following crates:

\begin{table}[h]
\centering
\small
\caption{Workspace crates, test counts, and brief descriptions (v0.1.2)}
\label{tab:status}
\begin{tabular}{p{2.5cm}lp{5.5cm}}
\toprule
Crate & Tests & Description \\
\midrule
\texttt{ternlang-core}    & 55  & Lexer, parser (+ \texttt{?} propagation), AST, semantic checker, BET VM (51 opcodes), ModuleResolver \\
\texttt{ternlang-ml}      & 21  & BitNet quantization, sparse/CSC matmul, deliberation engine, coalition vote, action gate, hallucination score \\
\texttt{ternlang-moe}     & 24  & MoE-13 orchestrator: dual-key routing, triad synthesis, 3-tier memory, 13 dual-signal agents, introspective hold \\
\texttt{ternlang-codegen} & 8   & AST-to-C11 transpiler; match/struct/propagation; forward declarations; native-compilable output \\
\texttt{ternlang-test}    & 10  & Full-pipeline test framework: \texttt{TernTestCase}, \texttt{assert\_tern!}, parse/semantic/VM error classification \\
\texttt{ternlang-hdl}     & 21  & Verilog-2001 codegen, BET RTL simulator (\texttt{BetRtlProcessor}), \texttt{ternlang sim --rtl} \\
\texttt{ternlang-runtime} & 4   & Distributed TCP actor runtime (\texttt{TernNode}) \\
\texttt{ternlang-compat}  & 29  & \texttt{.tasm} assembler (9-trit RISC), Owlet S-expr parser \\
\texttt{ternpkg}          & 5   & Package manager, GitHub-backed registry \\
\texttt{ternlang-lsp}     & --- & LSP 3.17: hover, completion (19 snippets), diagnostics \\
\texttt{ternlang-mcp}     & --- & MCP server, 10 tools: \texttt{trit\_decide}, \texttt{moe\_orchestrate}, \texttt{trit\_action\_gate}, \ldots \\
\texttt{ternlang-api}     & --- & REST API (18 endpoints), multi-tenant key management, SSE streaming \\
\bottomrule
\end{tabular}
\end{table}

The full test suite comprises \textbf{177+ tests across 12 crates}, all passing.
All crates are published to \texttt{crates.io} under open-core licensing
(LGPL-3.0-or-later for the compiler and CLI; BSL-1.1 converting to LGPL in
2030 for the ML engine and MoE orchestrator). The live API is deployed at
\texttt{https://ternlang.com/mcp} (Fly.io, Frankfurt) with TLS via Let's Encrypt,
serving both the marketing website and the MCP endpoint from a single Rust
binary. A companion \texttt{examples/} directory contains 265 executable
\texttt{.tern} programs spanning aerospace, medicine, distributed systems,
hardware pipelines, civic governance, finance, AI agent design, and Qutrit
Neural Network experiments—each catalogued in \texttt{INDEX.md} with patterns,
attribution, and cross-references. The repository is publicly available at the
URL in the author affiliation. Continuous integration via GitHub Actions
automatically builds, tests, and deploys to Fly.io on every push to main.

% ─────────────────────────────────────────────────────────────────────────────
\section{Conclusion and Future Work}
\label{sec:conclusion}

We have presented Ternlang v0.1.2, a full-stack balanced ternary execution
architecture spanning language design, ISA, virtual machine, hardware backend,
distributed runtime, native C compilation, and AI reasoning infrastructure. Five
layers of contribution are unified under the same foundational primitive—the
trit $t \in \{-1, 0, +1\}$ with an active neutral state:

\begin{enumerate}
  \item \textbf{Language completeness}: exhaustive three-way matching, the
        \texttt{?} error propagation operator, a real cross-file module system
        (stdlib built-in + filesystem-relative user modules), structured
        diagnostic codes, and all major control flow constructs—making ternlang
        a viable general-purpose programming language rather than a research
        prototype.

  \item \textbf{Execution substrate}: \texttt{TSPARSE\_MATMUL} achieves a
        $2.27\times$ reduction in multiply operations at baseline sparsity;
        a CSC kernel with Rayon reaches $122\times$ at high sparsity (512²
        matrix, 99\% sparsity, release mode). A native C11 compilation path
        via \texttt{ternlang-codegen} provides an alternative to VM
        interpretation.

  \item \textbf{Reasoning primitives}: the five-tool Ternary AI Reasoning
        Toolkit (DeliberationEngine, coalition vote, action gate, scalar
        temperature bridge, hallucination score) makes AI agent decision loops
        structurally three-valued.

  \item \textbf{MoE orchestration}: the MoE-13 orchestrator~\cite{moe13_2026}
        extends with 13-agent dual-signal deliberation, an introspective hold
        (stable attractor when $|\text{affirm} - \text{conflict}| \leq 1$
        across $\geq 4$ agents), and an \texttt{orchestrate\_full()} pipeline
        that pre-filters queries through the agent tier before MoE routing.

  \item \textbf{Deployment and tooling}: the live API at
        \texttt{https://ternlang.com/mcp} serves 10 MCP tools to any compatible
        AI client; a VS Code extension with LSP, formatter, and REPL provides
        an IDE experience; GitHub Actions CI/CD automatically builds, tests, and
        deploys on every push; and 177+ tests across 12 crates validate the
        full stack.
\end{enumerate}

Near-term publication targets include the VS Code Marketplace submission
(\texttt{ternlang-0.1.0.vsix}) and academic outreach to the USN group
(Bos \& Gundersen) for joint hardware-software co-design on the memristor
backend.

Further research directions include:

\begin{itemize}
  \item \textbf{Type inference}: Hindley-Milner style inference to eliminate
        explicit type annotations, reducing boilerplate for trit-returning
        functions
  \item \textbf{FPGA synthesis}: full \texttt{bet\_processor} targeting
        Xilinx Artix-7 and Lattice ECP5, benchmarked against the pure-Rust RTL
        simulator
  \item \textbf{End-to-end training}: native ternlang training loop with
        BitNet-style gradient quantization, ternary weight persistence, and
        backward pass in balanced ternary arithmetic
  \item \textbf{QNN acceleration}: formal integration of the Qutrit Neural
        Network framework~\cite{kepp2026qnn} with \texttt{TSPARSE\_MATMUL} as
        the inner-loop kernel, targeting quantum-adjacent hardware via the
        TRS stabilisation pattern
  \item \textbf{Memristor backend}: integration with physical ternary storage
        via the USN uMemristorToolbox~\cite{bos2020ternary}
  \item \textbf{Community bootstrap}: coordinated engagement with the broader
        ternary computing ecosystem (Brandon Smith 9-trit, Owlet S-expr,
        Trit-Rust, BitNet b1.58 community)
\end{itemize}

The ternary computing field has been fragmented for decades. Ternlang is
designed to be the substrate where those fragments converge—from ISA to
orchestrator, from single trit to 13-expert MoE, from quantum-adjacent
qutrit networks to physical memristor hardware. The philosophy is simple:
binary sees yes and no; ternary also sees \emph{not yet}—and that changes
everything.

% ─────────────────────────────────────────────────────────────────────────────
\bibliographystyle{IEEEtran}
\bibliography{references}

\end{document}
